\scanpage{101}
\mspage{50}
\deleted{et le diagramme multicartésien correspondant}
\deleted{\diagram{\begin{tikzcd}[ampersand replacement=\&]
X'_1\arrow[r]\arrow[d] \& X_1\arrow[r]\arrow[d] \& X'\arrow[r]\arrow[d] \& X=\unsure{/|1}\arrow[d] \\
 S^{\unsure{''|'}}\arrow[r,"f"] \& S\arrow[r,"g"] \& S'\arrow[r,"f"] \& S
\end{tikzcd}}}

La flèche $X'=f^*(X)\unsure{\longrightarrow|\text{---}}X$ s'interprète dans $\mathrm{Hot}_W$ comme
\diagram{\begin{tikzcd}[ampersand replacement=\&, column sep=small]
X'\arrow[rr]\arrow[d,no head] \& \& X\arrow[d,no head] \\
S'\arrow[rr,"f" description,no head]\arrow[dr,"p'"',no head] \& \& S\arrow[dl,"p",no head] \\
\& e \&
\end{tikzcd}}
\[
p'_!(f^*(\unsure{\xi|\zeta}))\longrightarrow p_!(\xi)\qquad\text{où }\xi=\mathrm{hot}_{S,W}(\unsure{X}),\tag*{(*)}
\]
Je dis que cette flèche est iso, i.e. que pour ts $y\in\mathrm{Hot}_W$, la flèche transposée
\[
\deleted{\unsure{\mathrm H}}\mathrm{Hom}_W(p_!(\xi),y)\longrightarrow\mathrm{Hom}(p'_!(f^*(\xi)),y)\tag*{(**)}
\]
est \unsure{bijection|bijective}. Admettant \unsure{la formule|les formules} d'adjonction, la flèche \unsure{(**)} s'interprète comme
\[
\mathrm{Hom}_W(\xi,\underbrace{p^*(y)}_{y_S})\longrightarrow\mathrm{Hom}_W(f^*(\xi),\underbrace{p'{}^*(y)}_{f^*(y_S)})
\]
i.e.
\[
\mathrm{Hom}_W(\xi,y_S)\longrightarrow\mathrm{Hom}_W(f^{\unsure{*|e}}(\xi),f^*(y_{\unsure{S|s}})),
\]
que ce soit iso résulte donc de $f^*$ pl. fid. sur les $\mathrm{HOT}_W^{\mathrm{lc}}$. Or on a vu que $f^*$ est une équivalence (cor. 2. page 46).
% archive sheet number: 101
